\documentclass[a4paper,14pt]{extarticle}
\usepackage{array,amsmath,fancybox,amsfonts,graphicx,amsthm,amssymb,wasysym,latexsym,slashed,anysize,subfigure,calc}
\usepackage[spanish]{babel}
\usepackage[utf8]{inputenc}
\usepackage{fancyhdr,float,hyperref}
\usepackage{lettrine}

\usepackage{enumitem}


%fractiones grandes en inline text: inventa un comando \ddfrac
\newcommand\ddfrac[2]{\frac{\displaystyle #1}{\displaystyle #2}}

%lewis y química
\usepackage{chemformula,elements}
%chemformula, provides \ch{} with the !(<below>)(<formula>) syntax and \chlewis;
%elements, provides \elconf and \writeelconf.
%ex:
%\ch{
%  !(\elconf{Li})( "\chlewis{180.}{Li}" ) +
%  !(\elconf{F})( "\chlewis{0.90:180:270:}{F}" )
%   ->
%  !(\writeelconf{2})( Li+ ) +
%  !(\writeelconf{2,2+6})( "\chlewis{0:90:180:270:}{F}" {}- )
%}


%entorno para química: texto y math mode
%\newcommand*\chem[1]{\textup{#1}}
\newcommand*\chem[1]{\ensuremath{\mathrm{#1}}}

%column
\usepackage{multicol}
%uso: \begin{multicols}{3}

%enumitem: configurar listas
\usepackage{enumitem}

%para los entornos de soluciones
\usepackage{mdframed,xcolor,color}
	%uso: 	entorno (begin,end): {mdframed}[backgroundcolor=brown!10] 

%subrayado y highlight. Tiene que estar detrás de usepackage[color]
\usepackage{soul}
	%\textst{normal} 
	\DeclareRobustCommand{\hlcyan}[1]{{\sethlcolor{cyan}\hl{#1}}} %en cyan!
		%\hl{highlighted text} %uso


%tamaños fuentes extra: 8,9,14,17,20
	\usepackage{extsizes}
	
 	%fuente librebaskerville
 	\usepackage[T1]{fontenc}
	\usepackage{librebaskerville}
	
	%fracciones $$ en texto más elaboradas
	 \usepackage{xfrac}
	 
	%permitir que una ecuación larga esté separada por un salto de página
	\allowdisplaybreaks[4]
		
	%gráficos varios 
	\usepackage{graphics}
	
	%fancy tables
	\usepackage{booktabs}
	\usepackage[small,centerlast,up]{caption}

	
	%mejora en la exposición de guarismos
	\usepackage{numprint}
	
\usepackage{everypage,draftwatermark} %Marca de agua
\SetWatermarkText{\vspace*{2cm}

\hspace*{2cm}\parbox{2\textwidth}{\centering DIST.\\B1-FyQ}}
\SetWatermarkScale{1.5}
\SetWatermarkLightness{0.9}

\DeclareMathOperator{\lin}{lin}
\DeclareMathOperator{\id}{\textbf{id}}
\DeclareMathOperator{\modulo}{mod}

\newcommand{\mc}[3]{\multicolumn{#1}{#2}{#3}}
\newcommand{\halmos}{
\vspace*{-4ex}

\begin{flushright}
\rule{1mm}{2.5mm} \end{flushright}
\vspace*{-3.5ex}

}
\renewcommand{\qed}{\halmos}


%versores: %uso: \uveci\ne\hat{\imath}_{\uvec{i}}+\uvec{j}+\uvec{k}$
\usepackage{bm}
\newcommand{\uveci}{{\bm{\hat{\textnormal{\bfseries\i}}}}}
\newcommand{\uvecj}{{\bm{\hat{\textnormal{\bfseries\j}}}}}
\DeclareRobustCommand{\uvec}[1]{{%
  \ifcsname uvec#1\endcsname
     \csname uvec#1\endcsname
   \else
    \bm{\hat{\mathbf{#1}}}%
   \fi
}}


%algún nuevo operador
\DeclareMathOperator{\cotan}{cotan}


%para textbf en entorno enumerate
\renewcommand{\labelenumi}{%
 \textbf{\theenumi}.
}

%a4
\marginsize{.5cm}{1cm}{1cm}{1cm}%requiere el paquete anysize,izq,dcha,arriba,abajo

%a5
%\marginsize{1cm}{1cm}{1cm}{.5cm}%requiere el paquete anysize,izq,dcha,arriba,abajo

\newtheorem{teo}{{\sc \textbf{Teorema}}}
\newtheorem{defi}{{\sc \textbf{Definición}}}
\newtheorem{coro}{{\sc \textbf{Corolario}}}


%opening
\title{\textbf{	}}
\date{}
\pagestyle{fancy}
\fancyhead[LO,RE]{{\small DOSSIER DE PROBLEMAS\hspace{1ex} \textbf{3Ev}}}
\fancyhead[RO,LE]{{\small B1-DIST\hspace{2ex} FÍSICA y QUÍMICA}}
%\fancyhead[CO,CE]{Ricardo Torres\\http://tinyurl.com/iesCMGfyq}



\begin{document}
\pagenumbering{gobble}

%gas ideal y mezclas
%disoluciones
%estequiometría: ajuste, reactivo limitante, reactivos en disolución e impuros


\begin{enumerate}
  \setlength\itemsep{1.5em}
    
    
    \item Se lanza, en vertical y hacia arriba, una pelota con velocidad $v_1=20\;ms^{-1}$. Desde la misma vertical y un segundo más tarde se hace lo propio con otra bola, esta vez con una velocidad inicial de $v_2=1.5v_1$. Despreciando rozamientos con el aire, determina:
    	\begin{itemize}
    		\item[a)] Las ecuaciones cinemáticas de las dos pelotas.
    		\item[b)] ¿Colisionan en algún punto? Si es así, determina cuándo y a qué altura.
    	\end{itemize}
    \qed 
    
    \item Formando un ángulo $\alpha=30^\circ$ con la horizontal se lanza un proyectil desde un cañón antiaéreo capaz de lanzar objetos con una velocidad de salida igual a $400\;ms^{-1}$. Sabiendo que la boca del cañón está a metro y medio del suelo, y despreciando cualquier rozamiento, determina: 
    	\begin{itemize}
    		\item[a)] Las ecuaciones cinemáticas vectoriales para el proyectil en cualquier punto de su vuelo.
    		\item[b)] ¿Cuál es la velocidad \textit{mínima} del proyectil en su vuelo?
    		\item[c)] ¿Cuánto tarda en colisionar con el suelo? ¿Cuál ha sido el alcance?
    		\item[d)] ¿Cuál es la altura máxima respecto al suelo que alcanza en su trayectoria?
    	\end{itemize}
    \qed
    
    \item Desde una altura $h$ se lanza un objeto con una velocidad $v_0$ formando un ángulo $\beta$ con la horizontal. Suponiendo que la intensidad gravitatoria es $g$, estima:
    	\begin{itemize}
    		\item[a)] El vector posición del objeto en cada instante de su vuelo.
    		\item[b)] La ecuación de la trayectoria del objeto\footnote{Se denomina trayectoria del objeto a la función $y(x)$ que expresa la componente $y$ del objeto en función de la componente $x$.}.
    		\item[b)] El ángulo, $\beta_{max}$, que maximiza el alcance. 
		\end{itemize}    	 
    \qed
    
    \item Los vectores de posición de dos barcos en el mar Mediterráneo vienen dados, en el {\small SI}, por las expresiones
    \[\vec{r}_1(t)=(t,t+e^{-t}),\hspace{2cm}\vec{r}_2(t)=\left(\sin t,1+t^{-1}+\cos t\right).\]
    	\begin{itemize}
    		\item[a)] Obtén $\vec{v}$ y $\vec{a}$  para cada embarcación.
    		\item[b)] Calcula la distancia que los separa en función del tiempo $t$.
    	\end{itemize}
    
\end{enumerate}

\end{document}

